\documentclass{article}
\usepackage[utf8]{inputenc}
\usepackage[margin=0.8in]{geometry}
\usepackage{graphicx}
\usepackage{pgffor}
\usepackage{caption}
\usepackage{paracol}
\usepackage{chngcntr}
\usepackage{amsmath}
\usepackage{biblatex}


\addbibresource{references.bib}

\renewcommand{\b}[1]{\mathbf{#1}}
\newcommand{\singleWidth}{0.6\textwidth}

\captionsetup[figure]{
    font=small,           % Makes text smaller (can use 'footnotesize' for even smaller)
    labelfont=bf,         % Makes "Figure X" bold
    margin=10pt,          % Indents caption from the left/right of the column
    skip=5pt,             % Reduces space between image and caption
    justification=centering % Keeps it neat
}


\title{Assignment 2 Report}
\author{Group 17: \\ Darlington Nkrumah (202492437) \\
  Greg de Souza (202582225) \\
  Xuan Toan Doan (202583882)}
\date{February 2026}

\begin{document}

\maketitle

\section*{Introduction}

For assignment 2 we were provided a two data sets, labeled ``train.csv''
and ``test.csv'' containing 8 features measuring some property of concrete,
and the value of the concrete compressive strength. For the sake of readability,
the 8 features will be labeled as $x_i$ according to the table \ref{tab:featvar},
and the target variable (Compressive Strength) will be $y$. Vectors
will be written with bold uppercase letters $\b{X}$ for a feature vector,
and $\b{\theta}$ for a set of parameters, with $\theta_0$ reserved for the intercept.
The collection of feature vectors from a
data set will be labeled $\mathcal{D}$.

In all questions the quality of fit was accessed with the $R^2$ metric, and the Err with the RSE metric

\begin{table}[h!]
  \centering
  \begin{tabular}{lc}
    \textbf{Feature Name}                              & \textbf{Variable} \\ \hline
   \textit{ Cement\_component1\_\_kgInAM\_3Mixture\_}           & $x_1$             \\
    \textit{BlastFurnaceSlag\_component2\_\_kgInAM\_3Mixture\_} & $x_2$             \\
    \textit{FlyAsh\_component3\_\_kgInAM\_3Mixture\_ }          & $x_3$             \\
    \textit{Water\_component4\_\_kgInAM\_3Mixture\_}            & $x_4$             \\
    \textit{Superplasticizer\_component5\_\_kgInAM\_3Mixture\_} & $x_5$             \\
    \textit{CoarseAggregate\_component6\_\_kgInAM\_3Mixture\_}  & $x_6$             \\
    \textit{FineAggregate\_component7\_\_kgInAM\_3Mixture\_}   & $x_7$             \\
    \textit{Age\_day\_}                                         & $x_8$             \\ \hline
    \textit{ConcreteCompressiveStrength\_MPa\_Megapascals\_}    & $y$                
  \end{tabular}
  \caption{Feature names with their symbolic variables}
  \label{tab:featvar}
\end{table}

  \section*{Question 1: Simple Linear Regression}

  The first question tasks us to train a predictor for $y$ using an ordinary least squares with a traditional
  linear regression model, such as in equation \ref{eq:lin_model}.
  The model was assessed for generalized error \textit{Err} in two different ways. First the training
  data set was split in training/validation/test, with 20\% of the data being randomly selected to be the
  test data set. Since there are no hyperparameters involved, the model was trained on the remaining
  ``Train + Validation'' dataset, consisting of 640 data points, and assessed for Generalized Error (Err)
  against the remaining 160 data points of the test set. The results are $R^2 = 0.6045$ and RSE = $11.5$,
  no variance is available since it's a single test dataset. The parameters fitted by this models are
  presented in the table \ref{tab:params}.

  \begin{equation}
    \label{eq:lin_model}
    \hat{y} = \theta_0 + \b{X}^T \b{\theta}
  \end{equation}

  The same model was also assessed with a cross validation approach with 5 folds
  over the ``train.csv'' dataset, meaning each folds had
  160 data points, similar to the test dataset from the previous approach. The choice of the number of folds
  seems adequate due to the fact that the number of remaining training data points  (640) still being
  way larger than the number of features (8). As an objective measure to validate this choice
  of number of folds, we have also evaluated the variance of the estimated Err in the cross validation.
  The results were $R^2 = 0.6270 \pm 0.0002$ and RSE $=11.6 \pm 0.3$, the low variance suggests that
  the 5 fold split results in representative values for the Err metrics. Also, that's the split that
  maximizes the $R^2$, as shown in figure \ref{fig:Q1}

  
  
  \begin{figure}[h!]
    \centering
    \includegraphics[width=0.8\textwidth]{Q1_figure1.png} 
    \caption{Assessment parameters given different number of folds for the Cross Validation Approach}
    \label{fig:Q1}
  \end{figure}


  The results for both the split and cross validation approaches is presented in table \ref{tab:Q1r2err},
  while the results suggests neither method found an adequate fit, the cross validation approach
  accesses the model as worse than the split approach. This makes sense, as the with the cross validation we
  access the model multiple times with different folds, making this result are more trustworthy estimation
  of the generalization error, specially because we can also estimate the variance of the model.

  In general, the bad performance is also expected in this ``naive'' approach, as the data features
  don't ``well behaved'' in a linear sense, as can be seen in figure \ref{fig:dataset}
  
  \begin{table}
    \centering
    \begin{tabular}{l|c|c}
      & \textbf{Split} & \textbf{CV} \\ \hline
      $R^2$ & $0.6045$ & $0.6270 \pm 0.0002 $ \\
      RSE & $11.5$ & $11.6 \pm 0.3 $
    \end{tabular}
    \caption{Results for the estimated quality of fit and Err using $R^2$ and $RSE$ in the Train/Test split and in the Cross Validation Aprroach}
    \label{tab:Q1r2err}
  \end{table}

  \section*{Question 2: Ridge Regression}

  The second task involves fitting the same data using Ridge Regressions (equation \ref{eq:ridge}) using
  the training dataset ``train.csv''.
  For this task we employed Grid Search using cross validation with 5 folds, the search determined
  that the best value for the
  hyperparameter is $\alpha=8111$, notably the data wasn't regularized. Figure \ref{Q2} presents
  the values of the RSE metric for different values of alpha, and there is a clear minimum around
  the found value.
  
  
  \begin{equation}
    \label{eq:ridge}
    \hat{y} = \theta_0 + \b{X}^T \b{\theta} + \alpha \cdot {||\b{\theta}||{^2_2}}
  \end{equation}

  
  \begin{figure}[h!]
    \centering
    \includegraphics[width=0.8\textwidth]{Q2_results.png} 
    \caption{Assessment of RSE metric for different values of alpha for the Ridge Regression model}
    \label{fig:Q2}
  \end{figure}

  The \textbf{cross validation} assessment estimated the Err parameters as $R^2 = 0.6271$ and RSE$=10.63$.

  The model was then retrained using $\alpha=8111$ on the total 800 data points of the training dataset,
  and accessed against the independent test dataset ``test.csv'', resulting in $R^2 = 0.5952$ and
  RSE$=10.00$.   As expected the performance against unseen data was worse than evaluated by the the
  cross validation grid search.
  
  For the same dataset and test, the ordinary least squares (OSL) result in a $R^2=0.5953$
  and RSE$= 10.00$. So, given the variance estimated in question 1, both regression methods have as
  similar performance. This is consistent with figure \ref{fig:Q2}, where a Ridge regression for $\alpha$
  close to 0 seems to perform similar to the best performing $\alpha$, it's reasonable to conclude
  that a ``naive'' linear relation between or targets and the features isn't enough to generate a
  good predictor.
  
  Notably, while the performance is similar, the actual parameters found by Ridge Regression are
  significantly smaller than those found by ordinary regression, which makes sense, since the job
  of the parameter $\alpha$ is two limit the growth of the L2 metric of the parameter set.
  the best set of parameters obtained by Ridge Regressions is. The values of the final parameters
  are presented in table \ref{tab:params}.

  \begin{table}[h]
    \centering
    \begin{tabular}{l|c|c|c}
      \textbf{Parameter} & \textbf{Ordinary Least Squares} & \textbf{Ridge Regression} & \textbf{Lasso Regression} \\ \hline
      Intercept & -44.19                 & 7.153           & 11.96            \\
      $\theta_1 $ & 0.1291                 & 0.1169           & 0.1145           \\
      $\theta_2 $ & 0.1158                 & 0.0994           & 0.0967           \\
      $\theta_3 $ & 0.1035                 & 0.0858           & 0.0816           \\
      $\theta_4 $ & -0.1328                & -0.2163          & -0.2213          \\
      $\theta_5 $ & 0.2056                 & 0.1092           & 0.1553           \\
      $\theta_6 $ & 0.0245                 & 0.0075           & 0.0065           \\
      $\theta_7 $ & 0.0296                 & 0.0126           & 0.0099           \\
      $\theta_8$  & 0.1253                 & 0.1172           & 0.1178          
    \end{tabular}
    \caption{Final parameters fitted in Q1, Q2 and Q3 for Ordinary, Ridge and Lasso regression respectively}
    \label{tab:params}
  \end{table}

  \section*{Question 3: Lasso Regression}
  The second task involves fitting the same data using Lasso Regressions (equation \ref{eq:ridge}) using
  the training dataset ``train.csv''.
  For this task we employed Grid Search using cross validation with 5 folds, the search determined
  that the best value for the
  hyperparameter is $\alpha=0.0001$, notably the data wasn't regularized. As figure \ref{fig:Q3} shows,
  there is a plateau of equivalent vales of $\alpha$ between 0.0001 and 1, so any value of the
  hyperparameter would be perform similar to another.  The \textbf{cross validation}
  assessment estimated the Err parameters as $R^2 = 0.6268$ and RSE$=10.63$. The

  \begin{figure}[h!]
    \centering
    \includegraphics[width=0.8\textwidth]{Q3_results.png} 
    \caption{Assessment of RSE metric for different values of alpha for the Lasso Regression model}
    \label{fig:Q3}
  \end{figure}
  
  \begin{equation}
    \label{eq:ridge}
    \hat{y} = \theta_0 + \b{X}^T \b{\theta} + \alpha \cdot {||\b{\theta}||{^2_2}}
  \end{equation}

  The model was then retrained using $\alpha=0.0001$ on the total 800 data points of the training dataset,
  and accessed against the independent test dataset ``test.csv'', resulting in $R^2 = 0.5953$ and
  RSE$=10.00$. As expected the performance against unseen data was worse than evaluated by the the
  cross validation grid search.

  Again, the performance is similar to the last two regression approaches, suggesting that the best
  performance for a linear relation between features and data is indeed around 0.6 for $R^2$.
  But the Lasso regression also resulted in some
  values of the parameters being very close to 0,  notably for $theta_6$ and $theta_7$, suggesting
  that the respective features $x_6$ and $x_7$ don't have meaningful linear relation with our target
  $y$.

  \begin{table}
    \centering
    \begin{tabular}{l|c|c|c|}
      \textbf{Err} \textbf{OLS} & \textbf{Ridge} & \textbf{Lasso} \\ \hline
      $R^2$ & 0.5953 & 0.5952 & 0.5953 \\
      RSE & 10.00 & 10.00 & 10.00
    \end{tabular}
    \caption{Results for the estimated quality of fit and Err using $R^2$ and $RSE$ in the Train/Test split and in the Cross Validation Aprroach}
    \label{tab:ErrCompar}
  \end{table}


  
  \begin{figure}[h!]
    \centering
    \includegraphics[width=0.8\textwidth]{dataset.png} 
    \caption{Features plotted against target}
    \label{fig:dataset}
  \end{figure}

  
\end{document}
